\chapter{Introducción}\label{cap:introduccion}

\section{Propósito y alcance}\label{sec:intro-proposito}

Este manual describe cómo instalar, configurar, desplegar, operar y mantener
Qadra a partir de su repositorio de código, y remite por ruta a la
documentación del repositorio cuando un tema requiere más detalle.

Se dirige a quien instala y opera el sistema (administrador de sistemas o
desarrollador con manejo de la línea de comandos de Linux, de PostgreSQL y de
las herramientas de Rust y Node.js) y a quien revisa el trabajo, por ejemplo
un sinodal, y necesita comprobar que el prototipo puede reconstruirse a partir
de sus fuentes. Quedan fuera del alcance el uso de la interfaz, que
corresponde al manual de usuario; la justificación del diseño, que corresponde
al reporte técnico, y el registro de despliegues concretos o del estado de una
instalación en particular.

\section{Qué es Qadra}\label{sec:intro-qadra}

Qadra es el \textbf{Prototipo de un sistema web para la gestión de
expedientes digitales y actividad procesal de un despacho legal, con
servicios de integridad y no repudio}, desarrollado como Trabajo Terminal
2026-B136 en la ESCOM del IPN. Se compone de un servidor escrito en Rust
(espacio de trabajo de Cargo en \archivo{crates/}), cuyo ejecutable
principal atiende la API HTTP y ofrece además una interfaz de línea de
comandos; de dos interfaces web en \archivo{web/apps/}, la del despacho
(\archivo{qadra}) y la del portal de la autoridad certificadora
(\archivo{ca}), que cuenta con un ejecutable propio en el mismo espacio de
trabajo; de una base de datos PostgreSQL y un servicio Redis, y de una
autoridad certificadora (CA) y una autoridad de sellado de tiempo (TSA)
internas, cuyos guiones están en \archivo{pki/}.

\importante{Qadra es un prototipo académico. Su CA y su TSA son internas: los
certificados, firmas y sellos de tiempo que produce sirven para demostrar los
servicios de integridad y no repudio dentro del prototipo, y este manual no
les atribuye validez jurídica ni equivalencia con servicios de certificación
acreditados.}

\section{Relación con el resto de la documentación}\label{sec:intro-relacion}

El manual es un documento LaTeX independiente (\archivo{manual-tecnico/}),
sin referencias cruzadas con los otros cuerpos de documentación del
repositorio (\cref{tab:intro-documentos}). El \cref{cap:mapa} indica qué
archivo de \archivo{docs/} consultar para cada tema.

\begin{table}[H]
  \centering
  \caption{Otros cuerpos de documentación del
    repositorio.}\label{tab:intro-documentos}
  \begin{tabularx}{\textwidth}{L{3.4cm} L{3.4cm} Y}
    \toprule
    \textbf{Documento} & \textbf{Ubicación} & \textbf{Qué responde} \\
    \midrule
    Reporte técnico
      & \archivo{latex/}
      & Por qué el sistema es como es: análisis, diseño, implementación y
        pruebas. Sus anexos B, C y D amplían capítulos de este manual. \\
    Manual de usuario
      & \archivo{manual-usuario/}
      & Cómo realiza cada rol sus tareas frente a la pantalla. \\
    Documentación del repositorio
      & \archivo{docs/}, \archivo{README.md}
      & El detalle de referencia: contratos HTTP, procedimientos de
        despliegue y recuperación, informes de verificación y decisiones de
        arquitectura. \\
    \bottomrule
  \end{tabularx}
\end{table}

\section{Cómo leer este manual}\label{sec:intro-lectura}

Los capítulos siguen el orden de una puesta en marcha, pero pueden consultarse
por separado (\cref{tab:intro-lectura}). Para reconstruir
el prototipo en un equipo propio bastan los
\cref{cap:requisitos,cap:instalacion}, en ese orden.

\begin{table}[H]
  \centering
  \caption{Capítulo que debe consultarse según la
    necesidad.}\label{tab:intro-lectura}
  \begin{tabularx}{\textwidth}{Y L{4.6cm}}
    \toprule
    \textbf{Si necesita\ldots} & \textbf{Consulte} \\
    \midrule
    Saber qué equipo y qué programas se requieren
      & \Cref{cap:requisitos} \\
    Entender los componentes y cómo se comunican
      & \Cref{cap:arquitectura} \\
    Levantar el sistema en su propio equipo
      & \Cref{cap:instalacion} \\
    Configurar el sistema
      & \Cref{cap:configuracion} \\
    Publicar una versión en un servidor
      & \Cref{cap:despliegue} \\
    Respaldar, restaurar y atender las tareas recurrentes del servicio
      & \Cref{cap:operacion} \\
    Administrar la CA y el sellado de tiempo internos
      & \Cref{cap:ca} \\
    Usar la interfaz de línea de comandos
      & \Cref{cap:cli} \\
    Consultar la API HTTP
      & \Cref{cap:api} \\
    Verificar un paquete de evidencia de forma independiente
      & \Cref{cap:verificacion} \\
    Dar mantenimiento al sistema o continuar su desarrollo
      & \Cref{cap:mantenimiento} \\
    Diagnosticar una falla
      & \Cref{cap:problemas} \\
    Localizar el documento de referencia de un tema
      & \Cref{cap:mapa} \\
    \bottomrule
  \end{tabularx}
\end{table}

\section{Convenciones}\label{sec:intro-convenciones}

Las rutas, los comandos, las variables de entorno y los nombres de binarios
aparecen en \texttt{tipografía fija}; las rutas son relativas a la raíz del
repositorio, salvo que se indique otra cosa. Los dominios, cuentas y
contraseñas de los ejemplos son ficticios. Los recuadros \textbf{Nota}
aportan contexto; los recuadros \textbf{Importante} advierten de límites del
prototipo y de acciones que no pueden deshacerse o que comprometen claves,
respaldos o datos.
